\documentclass[11pt]{article}

% ============================================================
% Milestone 1 -- Project Proposal & Literature Review
% Structure follows the course template (Sec. 6.1).
% ============================================================
\usepackage[margin=1in]{geometry}
\usepackage[T1]{fontenc}
\usepackage[utf8]{inputenc}
\usepackage{amsmath,amssymb}
\usepackage{graphicx}
\usepackage{booktabs}
\usepackage[round]{natbib}
\usepackage{hyperref}
\hypersetup{colorlinks=true, linkcolor=blue!50!black, citecolor=blue!50!black, urlcolor=blue!50!black}
\usepackage{enumitem}
\usepackage{xcolor}

\newcommand{\todo}[1]{\textcolor{red}{\textbf{[TODO: #1]}}}

\title{\LARGE Project Proposal \& Literature Review \\[4pt]
\large Blind Image Deblurring with a Compact Transformer}
\author{
  Zainab Ahmed (ID: 2501782) \and
  Rawda Abdelkhalek (ID: 2500020)
}
\date{Milestone 1 --- 2 November 2026 \\{} CSE 681 --- Digital Image Processing}

\begin{document}
\maketitle

\begin{abstract}
Motion blur degrades images captured while the camera or the scene is moving, and
removing it from a single image without knowing the blur kernel (blind deblurring) is
a hard inverse problem. We plan to study whether a small, parameter-budgeted
Transformer or hybrid CNN--Transformer can deblur images competitively with a classical
Wiener-filter baseline at low cost. We will use a subsampled version of the GoPro
deblurring dataset as the main benchmark, plus a small synthetically blurred DIV2K
subset as a sanity check for the Wiener baseline, where the blur kernel is known.
Methods will be compared on identical test images using PSNR, SSIM, and parameter count.
\end{abstract}

\section{Problem Statement \& Motivation}
\label{sec:motivation}
A blurred image can be modeled as a sharp image convolved with a blur kernel (the point
spread function, PSF), plus noise. If the PSF is known, classical deconvolution
methods such as Wiener filtering or Lucy--Richardson \citep{richardson1972,lucy1974}
can restore the image. In real photographs, however, the kernel is unknown and often
varies across the image, which makes the problem \emph{blind}. Deep networks learn the
mapping from blurred to sharp images directly from paired examples
\citep{nah2017,kupyn2019}, but the best-performing models can be large
\citep{kong2023,mao2024}. Our goal is to build a deliberately small model and measure
how it compares to the classical baseline in accuracy and in parameter count.

This connects to the course material on image restoration and the Wiener-filter
pipeline (Week 7, as referenced in the project catalog). \todo{Confirm the exact
Week-7 topic name from the course syllabus.}

\section{Related Work}
\label{sec:related}
\begin{table}[h]
\centering
\caption{Key references.}
\label{tab:refs}
\begin{tabular}{@{}p{0.30\linewidth}p{0.12\linewidth}p{0.45\linewidth}@{}}
\toprule
\textbf{Reference} & \textbf{Year} & \textbf{Relevance to this project} \\
\midrule
\citet{richardson1972}; \citet{lucy1974} & 1972; 1974 & Lucy--Richardson deconvolution: classical iterative baseline that assumes a known PSF. \\
\citet{nah2017} & 2017 & Introduced the GoPro dynamic-scene deblurring dataset (our main benchmark). \\
\citet{zamir2022} & 2022 & Efficient Transformer for high-resolution image restoration; a reference design for compact attention. \\
\citet{kupyn2019} & 2019 & DeblurGAN-v2; heavier reference point named in the project catalog. \\
\citet{kong2023} & 2023 & Frequency-domain Transformer for deblurring (recent, efficiency-oriented). \\
\citet{mao2024} & 2024 & Local frequency-domain Transformer for deblurring (recent). \\
\bottomrule
\end{tabular}
\end{table}

Classical deconvolution recovers a sharp image when the PSF is known
\citep{richardson1972,lucy1974}; its weakness for real photographs is that the kernel
must be supplied or estimated. \citet{nah2017} addressed real dynamic-scene blur with a
multi-scale convolutional network and released the GoPro pairs of blurred and sharp
frames, which later work uses as a standard benchmark; the authors note that blur
kernels are not provided because they are unknown. \citet{kupyn2019} is the heavier
learned reference point suggested by the project catalog.

Recent work applies Transformers to deblurring. \citet{kong2023} observe that attention
can be computed in the frequency domain with an element-wise product instead of a matrix
multiplication, and use this in the decoder of an asymmetric encoder--decoder network.
\citet{mao2024} compute channel-wise self-attention inside local frequency-domain
windows combined with an MLP gating mechanism. Both build on the line of efficient
restoration Transformers represented by \citet{zamir2022}. These papers report strong
accuracy; we have \emph{not} yet checked whether their models fit our compute budget, nor
measured them ourselves. \todo{Add at least one truly lightweight deblurring Transformer
from 2023--2026 (the catalog asks for 2--3 such papers) and verify its details before
submission.}

We do not claim a research gap beyond the project's stated objective: studying the
accuracy-versus-parameter-budget trade-off of a small model against a classical baseline.

\section{Dataset \& Baseline Plan}
\label{sec:data}
\textbf{Main dataset (GoPro).} The GoPro dataset of \citet{nah2017} contains 3214 blurred/sharp
pairs: 2103 for training (22 sequences) and 1111 for testing (11 sequences). Because the
official split is by sequence, train and test frames do not come from the same video.
\textbf{Sanity-check dataset (DIV2K).} A small subset of sharp DIV2K images will be blurred
synthetically with a known linear motion kernel, so that the Wiener filter can be checked
in the setting where it is expected to work.

\textbf{Subsampling plan (proposed, to be confirmed).} Uniformly take every $k$-th frame
from each sequence, holding out whole sequences for validation so that validation frames
are not near-duplicates of training frames. For example, $k=5$ would give about 420
training and 222 test pairs. The exact counts, selection method, and random seed will be
recorded in the repository (\texttt{data/README.md}) and in the report.

\begin{table}[h]
\centering
\caption{Dataset summary.}
\label{tab:data}
\begin{tabular}{@{}p{0.2\linewidth}p{0.42\linewidth}p{0.28\linewidth}@{}}
\toprule
\textbf{Dataset} & \textbf{Subset used} & \textbf{Link / license} \\
\midrule
GoPro (GOPRO\_Large) & \todo{final counts; official split is 2103 train / 1111 test pairs} & \todo{official URL and license} \\
DIV2K & \todo{e.g., 20--30 images, synthetic kernel, seed} & \todo{official URL and license} \\
\bottomrule
\end{tabular}
\end{table}

\textbf{Baseline.} Wiener-filter deconvolution (course Week 7 pipeline); Lucy--Richardson is an
optional second classical baseline. The Wiener filter needs a PSF. On DIV2K the PSF is
known. On GoPro it is not, so we will use a single assumed linear-motion kernel whose
length and angle are chosen on the validation set, and we will state clearly that this
is a weak, non-blind setting that gives the baseline no per-image kernel estimation.
\todo{Confirm that this kernel strategy is acceptable, and cite the course source for the Wiener filter.}

\section{Proposed Method}
\label{sec:method}
\textbf{Tentative plan (not yet confirmed).} The model takes a blurred RGB image and predicts
the sharp image (or a residual added to the input). A small convolutional stem and
encoder extract local features; a few Transformer blocks with an efficient (for example
channel-wise) attention operate on the reduced-resolution features to capture global
context; a convolutional decoder reconstructs the image. Compactness will be measured by
total parameter count (and optionally inference time). The target parameter budget is
\todo{fix a number after the first Colab test}.

\textbf{Evaluation protocol.} PSNR and SSIM on the same held-out test images for every method,
with identical preprocessing and a fixed random seed, plus parameter count. We make no
claim that the proposed model will outperform the baseline, and we make no novelty
claim from combining CNNs and Transformers; the final comparison will be reported as
measured. A parameter-budget ablation is planned for the final milestone, as suggested
by the project catalog.

\section{Feasibility \& Compute Plan}
\label{sec:feasibility}
Compute tier: GPU-light. The team has no strong local GPU, so training will run on Google
Colab (free tier or similar; the available GPU type is not guaranteed). The classical
baseline runs on CPU. Time per training run has not been measured yet.
\todo{Measure one short training run on Colab and state the expected time per experiment.}
Subsampling (Section~\ref{sec:data}) keeps the data small enough for this setting.

\subsection{Minimum Viable Deliverable (MVP)}
Wiener baseline numbers (PSNR/SSIM) on the GoPro subset and the DIV2K sanity check, and one
small model trained on the GoPro subset and evaluated on the same test images with PSNR,
SSIM, and parameter count, giving one valid comparison table.

\subsection{Stretch Goal}
If the MVP is on track by Weeks 5--6: the parameter-budget ablation with several model
sizes, Lucy--Richardson as a second baseline, and qualitative crops.
\todo{Decide whether a small robotics-motivated extension (e.g., effect on a
downstream task) is wanted; it is not included in the scope of this proposal.}

\section{Risk Assessment}
\label{sec:risk}
\begin{itemize}[leftmargin=*]
  \item \textbf{Unknown blur kernel for the Wiener baseline on GoPro.} Mitigation: DIV2K
        sanity check with known kernel; explicit statement of the assumed-kernel limitation.
  \item \textbf{Limited Colab compute and session time limits.} Mitigation: subsampled data,
        small model, checkpoints saved regularly, short runs first.
  \item \textbf{Team has limited deep-learning experience.} Mitigation: start from the smallest
        working model and verify each step with a short experiment before scaling up.
  \item \textbf{Dataset license or availability issues.} Mitigation: read the official license
        before use; fall back to the DIV2K synthetic variant if GoPro cannot be used.
  \item \textbf{Short milestone timeline.} Mitigation: the MVP is defined as the minimum
        valid comparison, and extras are optional.
\end{itemize}

\section{Work Breakdown}
\label{sec:workbreakdown}
The course template assumes three students; our team has two, so the three catalog roles
(Data \& Baseline, Model, Evaluation) are divided between two people.
\todo{Confirm the division below and confirm with the instructor that two-person teams are accepted.}

\begin{table}[h]
\centering
\caption{Milestone task allocation (draft; adapted from the project catalog entry).}
\label{tab:tasks}
\begin{tabular}{@{}p{0.12\linewidth}p{0.40\linewidth}p{0.40\linewidth}@{}}
\toprule
\textbf{Milestone} & \textbf{Zainab Ahmed} & \textbf{Rawda Abdelkhalek} \\
\midrule
M1 (Wk 3) & \todo{e.g., model design and evaluation protocol (PSNR/SSIM/params)} & \todo{e.g., GoPro/DIV2K subset pipeline and Wiener baseline} \\
M2 (Wk 6) & \todo{e.g., compact model and first trained checkpoint} & \todo{e.g., full Wiener baseline on the subset and preliminary comparison} \\
M3 (Wk 9) & \todo{e.g., final model and parameter-budget ablation} & \todo{e.g., final baseline table, full comparison, qualitative crops} \\
\bottomrule
\end{tabular}
\end{table}

\paragraph{AI-assistance disclosure.} Code scaffolding and drafting support for this
proposal were AI-assisted (Claude); all references were checked and all results will be
run and verified by the team.

\bibliographystyle{unsrtnat}
\bibliography{references}

\end{document}
